\chapter{Laplacian-eigenmode static sources}
\label{chap:laplace-sources}

\section{Laplace trial state}
The Laplace-static construction replaces the spatial Wilson line by a sum of Laplacian eigenvector outer products. With optional eigenmode weights \(\rho_i\), the source operator between \(\bm x\) and \(\bm y\) is built from
\begin{equation}
  \Phi(\bm x,\bm y;t)
  = \bar Q(\bm x,t)\left[\sum_{i=1}^{N_v}\rho_i^2 v_i(\bm x,t)v_i^\dagger(\bm y,t)\right]Q(\bm y,t).
  \label{eq:laplace-trial-state}
\end{equation}
The finite truncation in \(N_v\), or a non-trivial profile \(\rho_i\), is essential: if the complete eigenbasis is used with trivial weights, the source collapses to a local identity rather than an extended \QQbar\ operator.

\section{Static perambulators}
For a static color source at spatial position \(\bm x\), the temporal Wilson line from source time \(t_0\) to sink time \(t_1\) is projected into the Laplacian subspace through the static perambulator
\begin{equation}
  \tau_{ij}(\bm x;t_0,t_1)
  = v_i^\dagger(\bm x,t_0)\,U_t(\bm x;t_0,t_1)\,v_j(\bm x,t_1).
  \label{eq:static-perambulator}
\end{equation}
The resulting Laplace-static correlator is
\begin{equation}
  L(R,\tau)=\left\langle\sum_{i,j=1}^{N_v}\rho_i^2(t_0)\rho_j^2(t_1)
  \tau_{ij}(\bm y;t_0,t_1)\tau_{ji}(\bm x;t_1,t_0)\right\rangle,
  \label{eq:laplace-correlator}
\end{equation}
with \(R=|\bm y-\bm x|\) and \(\tau=t_1-t_0\). The static potential is then extracted from
\begin{equation}
  \Veff(R,\tau)=\log\frac{L(R,\tau)}{L(R,\tau+a)}.
  \label{eq:veff-laplace}
\end{equation}

\section{Profile choices and later extensions}
The current production tests use \(N_v=10\) with constant weights \(\rho_i=1\). This is intentionally simple: it validates the data flow and correlator normalization before introducing optimized Gaussian profiles or a GEVP basis. The literature indicates that larger eigenvector bases and Gaussian profile functions can improve ground-state overlap and lead to earlier plateaus \cite{Hoellwieser2023}.

\todo{Later update this chapter with the final profile strategy used in the thesis: constant profile only, larger \(N_v\), Gaussian profile basis, or GEVP-optimized profiles.}
